\documentclass[11pt]{article}
\usepackage[margin=0.92in]{geometry}
\usepackage{amsmath,amssymb,amsthm,mathtools}
\usepackage[hidelinks]{hyperref}

\newtheorem{theorem}{Theorem}
\newtheorem{remark}{Remark}
\newcommand{\A}{\mathcal A}
\newcommand{\dist}{\operatorname{dist}}

\title{Optimal Radial Equalization for Modular Frames with Central Lengths}
\author{Run fa\_banach\_001, agent\_lane\_16}
\date{2026-08-11}

\begin{document}
\maketitle

\section*{Source question and result}

Let $\A$ be a unital $C^*$-algebra and give the standard left module
$\A^d$ the inner product
\[
 \langle x,y\rangle=\sum_{k=1}^d x_k y_k^*.
\]
Problem 2.10 of Krishna, arXiv:2207.12799, asks for the closest
equal-inner-product modular frame to a given $\varepsilon$-nearly
equal-inner-product modular frame $\{\tau_j\}_{j=1}^n$, where
\begin{equation}\label{eq:near}
 (1-\varepsilon)\frac dn\,1
 \le p_j:=\langle\tau_j,\tau_j\rangle
 \le(1+\varepsilon)\frac dn\,1
 \qquad(1\le j\le n)
\end{equation}
and $0<\varepsilon<1$.  The squared modular distance is
\[
 \dist(\tau,\omega)^2
 =\left\|\sum_{j=1}^n
 \langle\tau_j-\omega_j,\tau_j-\omega_j\rangle\right\|.
\]

The theorem below completely answers the question when the squared lengths
$p_j$ are central, in particular for every commutative $C^*$-algebra.  For
arbitrary $\A$ it still gives an explicit equal-inner-product frame, exact
distance, and a dimensionally sharp quadratic-in-$\varepsilon$ bound.

\section*{Idea of the Proof}

Normalize each vector along its own left-module ray.  Functional calculus
makes the resulting length exactly $\sqrt{d/n}$ and controls the new frame
bounds.  To prove optimality, normalize any competitor as well.  The cross
inner product of two unit module vectors is a contraction.  When $p_j$ is
central, its square root can pass through the cross term, so the elementary
inequality $q+q^*\le2$ becomes a positive-element lower bound for each
competitor displacement.  The radial choice attains every one of these
lower bounds simultaneously.

\begin{theorem}[radial equalization and central-length optimality]
\label{thm:main}
Set $c=\sqrt{d/n}$ and define
\begin{equation}\label{eq:radial}
 \omega_j=c\,p_j^{-1/2}\tau_j\qquad(1\le j\le n).
\end{equation}
Then the following hold.
\begin{enumerate}
\item $\{\omega_j\}_{j=1}^n$ is an equal-inner-product modular frame.
If $\{\tau_j\}$ has frame bounds $a,b$, then $\{\omega_j\}$ has frame
bounds
\[
 \frac{a}{1+\varepsilon},\qquad \frac{b}{1-\varepsilon}.
\]
\item Its distance is exactly
\begin{equation}\label{eq:exact}
 \dist(\tau,\omega)^2
 =\left\|\sum_{j=1}^n(p_j^{1/2}-c1)^2\right\|,
\end{equation}
and consequently
\begin{equation}\label{eq:bound}
 \dist(\tau,\omega)^2
 \le d\bigl(1-\sqrt{1-\varepsilon}\bigr)^2
 \le d\varepsilon^2.
\end{equation}
\item If every $p_j$ belongs to the center $Z(\A)$, then
$\{\omega_j\}$ is globally closest among all equal-inner-product modular
frames.  Indeed, for every collection $\{\nu_j\}$ satisfying
$\langle\nu_j,\nu_j\rangle=c^2 1$,
\begin{equation}\label{eq:global}
 \dist(\tau,\nu)^2\ge \dist(\tau,\omega)^2.
\end{equation}
\end{enumerate}
\end{theorem}

\begin{proof}
By \eqref{eq:near}, each $p_j$ is positive and invertible.  With
$a_j=cp_j^{-1/2}$, the left-module convention gives
\[
 \langle\omega_j,\omega_j\rangle
 =a_jp_ja_j=c^2 1.
\]
Furthermore,
\[
 \frac1{1+\varepsilon}1\le a_j^2
 \le\frac1{1-\varepsilon}1.
\]
For every $x\in\A^d$,
\[
 \langle x,\omega_j\rangle\langle\omega_j,x\rangle
 =\langle x,\tau_j\rangle a_j^2\langle\tau_j,x\rangle.
\]
Sandwiching the displayed order bounds and then using the original frame
inequalities proves the new frame bounds.

Since $1-cp_j^{-1/2}$ is a self-adjoint function of $p_j$,
\[
 \langle\tau_j-\omega_j,\tau_j-\omega_j\rangle
 =(1-cp_j^{-1/2})p_j(1-cp_j^{-1/2})
 =(p_j^{1/2}-c1)^2,
\]
which proves \eqref{eq:exact}.  The spectrum of $p_j^{1/2}$ lies in
\[
 \bigl[c\sqrt{1-\varepsilon},c\sqrt{1+\varepsilon}\bigr].
\]
For $0<\varepsilon<1$ the larger endpoint deviation is
$c(1-\sqrt{1-\varepsilon})$.  Hence each summand in \eqref{eq:exact} is at
most $c^2(1-\sqrt{1-\varepsilon})^2 1$.  Summing, using $nc^2=d$, and
observing
\[
 1-\sqrt{1-\varepsilon}
 =\frac{\varepsilon}{1+\sqrt{1-\varepsilon}}\le\varepsilon
\]
gives \eqref{eq:bound}.

It remains to prove optimality under centrality.  Put
\[
 v_j=p_j^{-1/2}\tau_j,\qquad u_j=c^{-1}\nu_j,
 \qquad q_j=\langle v_j,u_j\rangle.
\]
Then $\langle v_j,v_j\rangle=\langle u_j,u_j\rangle=1$.  The
$C^*$-module Cauchy--Schwarz inequality gives $\|q_j\|\le1$, whence
$q_j+q_j^*\le2 1$.  Since $p_j^{1/2}$ is central, expansion in the
left-module convention yields
\begin{align*}
 \langle\tau_j-\nu_j,\tau_j-\nu_j\rangle
 &=p_j+c^2 1-cp_j^{1/2}(q_j+q_j^*)\\
 &\ge p_j+c^2 1-2cp_j^{1/2}
  =(p_j^{1/2}-c1)^2.
\end{align*}
Summing preserves order.  Norm is monotone on positive elements, so
\eqref{eq:global} follows from \eqref{eq:exact}.  The radial collection
attains equality and is a frame by the first part.
\end{proof}

\section*{Scope and obstruction}

The centrality hypothesis is automatic when $\A$ is commutative, but it is
not cosmetic in the proof.  Without it, the two cross terms are
\[
 cp_j^{1/2}q_j\quad\text{and}\quad cq_j^*p_j^{1/2};
\]
they cannot be combined as positive multiplication of $q_j+q_j^*\le2$.
Composing with a state does not repair this, because a general state is not
tracial.  Likewise the modular distance is a $C^*$-norm of an
$\A$-valued sum, not a scalar unitarily invariant Hilbert--Schmidt norm.
Thus Theorem \ref{thm:main} is a full solution of Problem 2.10 for central
squared lengths and a controlled feasible solution in general; the optimal
noncentral frame remains open.

Eight focused upgrade routes were pursued: duplicate checking against the
paper's appendix, arbitrary-algebra radial normalization, exact frame and
distance estimates, a commutative fiberwise proof, the central-order
upgrade, statewise lower bounds, polar-decomposition analogies, and
alternating length/frame normalization.  The last three all meet the same
opposite-sided cross-term obstruction.

\section*{Novelty check}

Exact arXiv-id, title, ``closest equal inner product modular frame,'' and
``Modular Paulsen Problem'' searches were run in the local FA/Banach indexes
and primary arXiv sources through 2026-08-11.  No later answer to Problem
2.10 and no existing run packet for it were found.  The appendix
restricted-invertibility and literally stated Johnson--Lindenstrauss
questions are already handled by the separate GPT5.6 packet attached to
arXiv:2208.05223 and are not duplicated here.

\section*{References}

\begin{itemize}
\item K. Mahesh Krishna, \emph{Modular Paulsen Problem and Modular
Projection Problem},\\ arXiv:2207.12799.
\item M. Frank and D. R. Larson, \emph{Frames in Hilbert $C^*$-modules and
$C^*$-algebras}, J. Operator Theory 48 (2002), 273--314.
\end{itemize}

\end{document}
